\begin{exo}
	On considère deux événements $A$ et $B$ d'une même situation aléatoire tels que :
	\[
		P(A) = 0,3 \quad ; \quad P(B) = 0,6 \quad \text{et} \quad P(A \cap B) = 0,2
	\]

	\begin{enumerate}
		\item Rappeler la formule du cours qui permet de calculer la probabilité conditionnelle $P_{A}\pa{B}$ à l'aide de $P\pa{A}$ et de $P\pa{A \cap B}$.
		\item Calculer la probabilité que l'événement $B$ se réalise sachant que l'événement $A$ est réalisé.
		\item Calculer de la même manière la probabilité conditionnelle $P_{B}\pa{A}$.
	\end{enumerate}
\end{exo}
